\chapter{Carga del Firmware}

Este capítulo explica cómo instalar el entorno de desarrollo y cargar el firmware en el microcontrolador ESP32.
El firmware habilita la comunicación Bluetooth con la aplicación móvil y proporciona la lógica de control para
motores, servos, sensores y telemetría.

El procedimiento utiliza el IDE de Arduino, que ofrece un entorno sencillo para compilar y cargar firmware en
placas de desarrollo ESP32.

\section{Instalar el IDE de Arduino}

Descarga e instala el IDE de Arduino desde el sitio web oficial:

\begin{center}
    \url{https://www.arduino.cc/en/software}
\end{center}

Se admiten tanto Arduino IDE 1.8.x como Arduino IDE 2.x, pero se recomienda Arduino IDE 2.x porque ofrece mejor
rendimiento y una interfaz de usuario moderna.

Después de instalar el software, abre el IDE de Arduino para continuar con la configuración del ESP32.

\section{Instalar el Soporte de Placas ESP32}

El IDE de Arduino requiere el paquete de placas ESP32 para poder compilar firmware para microcontroladores ESP32.

Abre el IDE de Arduino y navega a:

\begin{verbatim}
    File -> Preferences
\end{verbatim}

Ubica el campo \textbf{Additional Boards Manager URLs} y agrega la siguiente dirección:

\begin{verbatim}
    https://dl.espressif.com/dl/package_esp32_index.json
\end{verbatim}

Presiona \textbf{OK} para guardar la configuración.

A continuación, instala las definiciones de la placa ESP32:

\begin{verbatim}
    Tools -> Board -> Boards Manager
\end{verbatim}

Busca:

\begin{verbatim}
    ESP32 by Espressif Systems
\end{verbatim}

Instala el paquete de la placa.

Para obtener la mejor compatibilidad con este proyecto, se recomienda la siguiente versión:

\begin{verbatim}
    ESP32 Core Version 2.0.x
\end{verbatim}

Esta versión proporciona soporte Bluetooth estable y funciona bien con la arquitectura del firmware utilizada en este proyecto.

\section{Descargar el Firmware}

Descarga el código fuente del proyecto desde el repositorio de GitHub:

\begin{center}
    \url{https://github.com/MICSBOL/ESP32_BT_Controller}
\end{center}

Puedes descargar el proyecto haciendo clic en el botón \textbf{Code} y seleccionando \textbf{Download ZIP},
o clonando el repositorio usando Git.

Después de descargar el proyecto, extrae los archivos en una carpeta de tu computadora.

\section{Abrir el Proyecto del Firmware}

Abre el IDE de Arduino y selecciona:

\begin{verbatim}
    File -> Open
\end{verbatim}

Navega hasta la carpeta que contiene el firmware y abre el archivo principal del proyecto:

\begin{verbatim}
    ESP32_BT_Controller.ino
\end{verbatim}

Cuando el archivo se abre, el IDE de Arduino cargará automáticamente todos los archivos fuente adicionales que usa
el firmware, incluyendo los componentes modulares responsables de la lógica de control, comunicación, telemetría
e interfaces de hardware.

\section{Seleccionar la Placa ESP32}

Configura el tipo de placa usado para la compilación.

Navega a:

\begin{verbatim}
    Tools -> Board
\end{verbatim}

Selecciona:

\begin{verbatim}
    ESP32 Dev Module
\end{verbatim}

Esta configuración es compatible con la mayoría de placas de desarrollo ESP32 disponibles en el mercado.

\section{Configurar Parámetros de la Placa}

Para un funcionamiento confiable, se recomienda usar la siguiente configuración de placa:

\begin{itemize}

    \item Board: ESP32 Dev Module

    \item Upload Speed: 921600

    \item CPU Frequency: 240 MHz

    \item Flash Frequency: 80 MHz

    \item Flash Mode: QIO

    \item Flash Size: 4MB

    \item Partition Scheme: Default 4MB

    \item Core Debug Level: None

    \item PSRAM: Disabled

\end{itemize}

Estos ajustes aseguran una compilación estable y una carga de firmware confiable para la mayoría de placas ESP32.

\section{Conectar la Placa ESP32}

Conecta la placa de desarrollo ESP32 a tu computadora usando un cable USB.

Después de conectar la placa, selecciona el puerto serie utilizado por el dispositivo:

\begin{verbatim}
    Tools -> Port
\end{verbatim}

Nombres típicos de puertos incluyen:

\begin{verbatim}
    COM3
    COM4
    /dev/ttyUSB0
    /dev/ttyACM0
\end{verbatim}

Selecciona el puerto correspondiente a tu dispositivo ESP32.

\section{Compilar el Firmware}

Antes de cargar el firmware, se recomienda verificar que el código compila correctamente.

Haz clic en el botón \textbf{Verify} en el IDE de Arduino.

El IDE compilará todos los archivos fuente usados por el proyecto. Si la compilación termina correctamente,
la consola mostrará un mensaje indicando el uso de memoria del programa.

\section{Cargar el Firmware}

Haz clic en el botón \textbf{Upload} para transferir el firmware compilado al ESP32.

Durante el proceso de carga, el IDE compilará el firmware y lo enviará al microcontrolador a través de la conexión USB.

Algunas placas ESP32 requieren presionar el botón \textbf{BOOT} cuando comienza la carga. Si la carga falla,
mantén presionado el botón BOOT mientras el firmware se está cargando.

Cuando la carga finaliza, el IDE mostrará un mensaje indicando que el ESP32 se está reiniciando.

\section{Verificar el Funcionamiento del Firmware}

Abre el Monitor Serie para confirmar que el firmware está ejecutándose.

Navega a:

\begin{verbatim}
    Tools -> Serial Monitor
\end{verbatim}

Configura la velocidad de comunicación serie en:

\begin{verbatim}
    115200
\end{verbatim}

Cuando el ESP32 arranca, deberías ver mensajes de estado impresos en la consola indicando que el controlador
Bluetooth se ha iniciado y está listo para aceptar conexiones desde la aplicación móvil.

\section{Siguiente Paso}

Después de cargar el firmware correctamente, el siguiente paso es conectar la aplicación móvil al ESP32 y comenzar
a controlar el dispositivo.

El proceso de conexión se describe en el Capítulo 5.